\honest{Hold Off On Proposing Solutions}{Eliezer Yudkowsky}{2007}

I quote Robyn Dawes at length. Norman Maier noticed that group members propose solutions as soon as they start discussing a problem, become ``emotionally attached'' to them, and never suggest better ones. His rule: ``Do not propose solutions until the problem has been discussed as thoroughly as possible without suggesting any.'' In his role-play, three workers and a foreman discuss an expert's advice to stop rotating jobs. Groups given the rule more often find the better answer: the two abler workers rotate and the least able keeps the easiest job. Dawes uses the rule with his own groups, especially on tough problems, though he has ``no objective criterion'' for judging the results, so the rule ``appears'' to help. \nb{Dawes's book and Maier's papers could not be checked here. The quoted passage gives no figures for the size of the effect.}

My verdict: ``This is so true it's not even funny.'' It gets worse as problems get harder. Many people I meet seem to know how to build an artificial general intelligence, and ``an actual majority'' solve Friendly AI within fifteen seconds. \nb{The footnote cites Yudkowsky's own chapter, which reports the same kind of encounter and gives no count.} Physicists, economists and evolutionary biologists meet the same thing.

Maier's rule echoes ``The Bottom Line'': the effectiveness of a decision is ``determined only by'' the thinking done in first reaching it. We change our minds less often than we think: 24 people gave their expected choice an average probability of 66\%, and only one chose otherwise. ``Once you can guess what your answer will be, you have probably already decided,'' and if you can guess it in half a second, you have half a second to be intelligent. \nb{The half second has no source; see ``We Change Our Minds Less Often Than We Think.''}

``Traditional Rationality'' stresses giving up an opinion in the face of clear evidence. But an idea, once in your head, will probably need too much evidence to remove, and strong evidence is not always available.

So, I suspect, a more powerful method is to ``hold off on thinking of an answer,'' and stretch out the moment before we can guess it. Even half a minute would be better than half a second. \nb{The post's one experiment concerned groups. John Dewey, in \textsc{How We Think} (1910), had already called suspended judgment ``the essence of critical thinking'': inquiry into the nature of the problem ``before proceeding to attempts at its solution.''}
